\section{Differentiable acquisition design}

\begin{figure}[t]
  \centering
  \includegraphics[width=\linewidth]{figures/differentiable-oed-double-loop-draft_ipnut_output.pdf}
  \caption{Nested differentiable OED. The inputs to the outer design loop are the initial receiver positions and initial velocity model. At outer design iteration $t$, measurements at
$\vect{p}^t$ drive $K$ FWI updates; the design objective then updates the
receivers to $\vect{p}^{t+1}$, where data are reacquired. The outputs of this workflow are the updated receiver positions and the corresponding reconstructed models. Red crosses denote fixed sources and yellow stars denote movable
  receivers.}
  \label{fig:diff-oed-workflow}
\end{figure}

\subsection{Full-waveform inversion}

Figure~\ref{fig:diff-oed-workflow} shows the bilevel structure of the workflow. The two optimization levels are indexed separately: $t$
indexes iterations of the outer design loop, and $k$ indexes FWI updates. Let
$\vect{m}$ denote the discretized velocity model,
$\vect{p}\in\mathbb{R}^{2N_r}$ the continuous receiver coordinates of $N_r$ receivers,
and $\mathcal{F}_{\vect{p}}(\vect{m})$ the time-domain acoustic forward
operator sampled at those coordinates. Let $\mathcal{U}_{m}$ denote the FWI
optimizer. For predicted data
$\vect{d}^{\mathrm{pred}}=\mathcal{F}_{\vect{p}}(\vect{m})$, the data misfit is
the mean squared error between observed and predicted waveforms,
\begin{equation}\label{eq:fwi_loss}
  \Ldata(\vect{d}^{\mathrm{pred}},\vect{d}^{\mathrm{obs}})
  = \left\|\vect{d}^{\mathrm{pred}}-\vect{d}^{\mathrm{obs}}\right\|_2^2.
\end{equation}

At outer design iteration $t$, the geometry is
$\vect{p}^{t}$ and the corresponding observed waveforms are
$\vect{d}^{\mathrm{obs}}(\vect{p}^{t})$. Starting from $\vect{m}_0^{t}$, $K$
steps of FWI produce
\begin{align}
  \vect{m}_{k+1}^{t}
  &= \mathcal{U}_{m}\!\left(
     \vect{m}_k^{t},
     \nabla_{\vect{m}}\Ldata\!\left(
       \mathcal{F}_{\vect{p}^{t}}\!\left(
         \vect{m}_k^{t}
       \right),
       \vect{d}^{\mathrm{obs}}(\vect{p}^{t})
     \right)
  \right),
  \qquad k=0,\ldots,K-1,
  \label{eq:fwi-update}
\end{align}
The latest velocity model is carried over to the next outer design iteration,
such that $\vect{m}_0^{t+1}=\vect{m}_K^{t}$. The complete workflow is differentiable and uses automatic differentiation to compute the FWI gradients as well as the outer design gradients.

\subsection{Interchangeable design objectives}

Here, we present the three design objectives used in this study. Each objective
uses the current receiver coordinates $\vect{p}^{t}$ and the corresponding
reconstruction state defined below. The objectives are interchangeable, and the
same workflow can be used to compare their performance.

\paragraph{Oracle model error.}
For synthetic validation, the privileged objective directly measures
reconstruction quality against the true model $\vect{m}^{\star}$,
\begin{equation}
  \mathcal{L}_{\mathrm{model}}(\vect{p}^{t})
  = \left\|\vect{m}_K^{t}-\vect{m}^{\star}\right\|_2^2.
\end{equation}
It provides an upper benchmark but is unavailable in a physical experiment.

\paragraph{Waveform misfit.}
Let $\vect{d}^{\mathrm{obs}}(\vect{p})$ denote measurements acquired after the
receivers have moved to $\vect{p}$. We evaluate the reconstructed model using
the same acquisition geometry:
\begin{equation}
  \mathcal{L}_{\mathrm{wave}}(\vect{p}^{t})
  = \left\|\mathcal{F}_{\vect{p}^{t}}
  \!\left(\vect{m}_K^{t}\right)
  -\vect{d}^{\mathrm{obs}}(\vect{p}^{t})\right\|_2^2.
  \label{eq:waveform}
\end{equation}
This objective does not require the unknown true material model. Because the
waveform misfit is evaluated at a different acquisition geometry at every
outer design iteration, its values across design iterations are not expected to
decrease monotonically.

\paragraph{Stochastic illumination.}
The illumination objective scores a receiver geometry by how strongly and how
uniformly the simulated data respond to perturbations throughout the model.
It uses the diagonal of the local Gauss--Newton matrix $J^\top J$, where
$J=\partial\mathcal{F}_{\vect{p}}(\vect{m})/\partial\vect{m}$ is evaluated at
the current reconstructed model. Large diagonal entries indicate model
locations to which the acquisition is locally sensitive. We estimate this
diagonal with random probes and vector--Jacobian products, so the full Jacobian
and Gauss--Newton matrix are never formed. The loss penalizes spatial
nonuniformity while rewarding a large mean sensitivity. This criterion requires neither the true model nor observed
waveforms. Appendix~\ref{app:illumination-derivation} derives the estimator and
gives the exact scalarization; its numerical settings are listed in
Table~\ref{tab:outer-design-parameters}.

\subsection{Design update and differentiation paths}

Let $\Ldesign$ denote one of the three objectives above and $\mathcal{U}_{p}$
the coordinate optimizer. The receiver-coordinate update is
\begin{equation}
  \vect{p}^{t+1}
  = \mathcal{U}_{p}\!\left(
       \vect{p}^{t},
       \nabla_{\vect{p}}\Ldesign(\vect{p}^{t})
     \right),
  \label{eq:design-update}
\end{equation}
The objectives induce two different differentiation paths. For the oracle
model error and waveform misfit
objectives, reverse-mode automatic differentiation propagates the
receiver-coordinate gradient through the complete unrolled FWI algorithm,
including all $K$ FWI updates and their wave simulations. By contrast, the
illumination gradient is differentiated through the forward wave solver and
its vector--Jacobian products, but not through the FWI trajectory.  Illumination therefore reduces the additional graph required to
calculate the design update. All methods
still perform $K$ FWI updates to advance the reconstruction between design
iterations. We implement the workflow in JAX using a differentiable wave solver
in the spirit of j-Wave~\citep{stanziola2023jwave}.

We do not assume that the inner FWI problem is fully converged. Instead, the
$K$ updates in Eq.~\eqref{eq:fwi-update} define the reconstruction presented to
the outer objective, so the design is optimized for a given number of FWI
updates. We fix $K$ in this study. In future, the inversion budget $K$, or a
learned stopping rule based on the level of convergence, could itself be
optimized. In this finite-iteration formulation, the chosen optimizer and
stopping point are hyperparameters, and retaining the updates determines the
memory cost of differentiation.
